\section{Source-only Scalar Calculation Ledger}
\label{sec:strategy2-pure-theorems}

This uncompiled source records the long scalar calculations associated with
Appendix~\ref{sec:strategy2-optimization-problems}.  It is a calculation
ledger, not the proof architecture used in the paper.  Except for the
explicitly bridge-realized T3 statement below, the labels with prefix
\texttt{thm:s2-univ-} denote universal real-variable assertions.  The active
paper uses distinct \texttt{thm:s2-geom-} labels for geometric consequences.
Nothing in this file formalizes a geometry-to-parameter bridge, a global case
split, or the final covering contradiction; Lean currently elaborates only
the corresponding scalar statement shells with intentional admissions.

\subsection{Primitive high-radial output bound}

\begin{lemma}[The pure cap inequality]
\label{lem:s2-pure-cap}
For every $1/2<c<1$ and $0\le x\le1$,
\[
 0\le F(c,x)\le1-x.
\]
The upper inequality is strict on every non-linear branch.
\end{lemma}

\begin{proof}
The lower bound is immediate from Definition~\ref{def:s2-pure-map}.  On a
linear branch the upper bound is equality.  On a $Q_-$ branch, where $x<c$,
\[
 Q_-(x,c)\le1-x
 \iff
 \sqrt{x^2-cx+c^2}\le c
 \iff x(x-c)\le0.
\]
On a $Q_+$ branch, $x<1$ and Lemma~\ref{lem:s2-pure-map-well-defined} gives
\[
 \sqrt{D_+(c,x)}\ge c(2c-1).
\]
After multiplying by the positive denominator in $Q_+$ and dividing by
$1-x$, this is exactly $Q_+(x,c)\le1-x$; equality can occur only at
$x=1-c$, which is assigned to linear.  On the low-$c$ tie,
$h(c)\le1/2=1-1/2\le1-h(c)$; equivalently
$\sqrt{16c^2-3}\ge4c-1$, whose square is $8c\ge4$.  On the high-$c$
plateau, $\ell(c)<x<r(c)=c-\ell(c)$ gives
$\ell(c)<1-c+\ell(c)<1-x$.  These cases exhaust the definition of $F$.
\end{proof}

\subsection{The four all-Vd0 problems}

\begin{theorem}[Pure S2-E1]
\label{thm:s2-univ-e1}
For every $(r,a,d)\in\mathcal D_1\cup\mathcal D_2$,
\[
 \Psi_{\rm E1}(r,a,d)<0.
\]
\end{theorem}

\begin{proof}
Write $w=1-r$, $S=e$, and
\[
 s=\frac{k}{r},\qquad u=r-d,\qquad
 \omega=1-s-u=w+d-\frac{k}{r}=\frac{\Delta_R}{r}>0.
\]
Then $0<u<r$, $s,u,\omega>0$, and $s+u+\omega=1$.  Put
$\delta_0=r/(1+S)$.  Using $S^2=1-rw$ and
$p=S(1-S)$, direct simplification gives
\[
 u-\delta_0-\frac r w\omega=\frac a w>0.
\]
Thus the variables $(R,S,s,u,w)$ in
Lemma~\ref{lem:one-side-capped-loss} may be taken as
$(r,e,s,u,\omega)$, and its two reach lower bounds are
\[
 c_1=\frac ur=1-\frac dr,
 \qquad
 c_5=1-\frac a w.
\]
Since the technical notation in that lemma satisfies
$\overline M_c(x)=F(c,x)$, its conclusion is precisely
$\Psi_{\rm E1}(r,a,d)<0$.
\end{proof}

\begin{lemma}[Pure paired-endpoint subcase]
\label{lem:s2-pure-paired-kernel}
Let $0<w\le1/2$, $r=1-w$, $K=\sqrt{1-wr}$, and put
\[
 z=\frac{w}{1+K},\qquad \zeta=\frac{r}{1+K}.
\]
Suppose $1/2<\alpha,\beta<1$, $p=w\alpha$, $q=r\beta$, and
\[
 \beta+p>1+z,
 \qquad
 \alpha+q>1+\zeta.
\]
Then
\[
 F(\alpha,p)+F(\beta,q)<1.
\]
\end{lemma}

\begin{proof}
The proof of Lemma~\ref{lem:two-endpoint-capped-loss}, beginning immediately
after the derivation of \eqref{eq:ce2-strict-endpoints}, is an algebraic
proof under exactly the hypotheses displayed here.  It first partitions the
two evaluations of $F$ into the four exact algebraic branches, obtains the seven
possible ordered pairs in \eqref{eq:two-endpoint-pairs}, and proves the strict
loss for each pair.  Every later quantity in that proof is an explicit
function of $(w,r,K,z,\zeta,\alpha,\beta,p,q)$; no placement, trace, or
existence assertion is invoked.  Hence the same branchwise calculation proves
the universal statement above.
\end{proof}

\begin{theorem}[Pure S2-E2]
\label{thm:s2-univ-e2}
For every $(r,a,d)\in\mathcal D_2$,
\[
 \Psi_{\rm E2}(r,a,d)<0.
\]
\end{theorem}

\begin{proof}
First assume $w\le r$.  Set
\[
 \alpha=\frac{w-a}{w},\quad p=w-a,
 \qquad
 \beta=\frac{r-d}{r},\quad q=r-d.
\]
Proposition~\ref{prop:s2-pure-domain-bounds} gives
$1/2<\alpha,\beta<1$.  Since $e^2=1-rw$ and
$1-e=rw/(1+e)$, the two strict subcase inequalities are equivalent to
\[
 \beta+p>1+\frac{w}{1+e}
 \iff \Delta_L>0,
\]
\[
 \alpha+q>1+\frac{r}{1+e}
 \iff \Delta_R>0.
\]
Lemma~\ref{lem:s2-pure-paired-kernel} therefore applies and gives the claimed
objective sign.  If $r<w$, interchange
$(r,d)$ with $(w,a)$.  The signed CE2 domain and the sum defining
$\Psi_{\rm E2}$ are invariant under this interchange, reducing to the first
case.
\end{proof}

\begin{theorem}[Pure S2-R1]
\label{thm:s2-univ-r1}
For every $(r,a,d)\in\mathcal D_1$,
\[
 \Psi_{\rm R1}(r,a,d)<0.
\]
\end{theorem}

\begin{proof}
On $\mathcal D_1$ one has
\[
 d+ra\ge p,\qquad a+wd<p,
 \qquad rd-wa>0,
\]
so $m=a/r$.  Substitute
\[
 R=r,\,W=w,\,A=a,\,D=d,\,
 X=r-d,\,h_0=\frac{k}{2r}
\]
into the scalar calculation in the proof of
Proposition~\ref{prop:ce1-one-gap}.  From
\eqref{eq:ce1-normal-domain} through the proof of
\eqref{eq:ce1-suffix}, every step uses only the displayed real inequalities,
the identity $e^2=1-rw$, and the explicit branches of $F$ and $T$.
The calculation proves
\[
 [T(1-a,\cdot)\mid T(1-m,\cdot)\mid T(1-d,\cdot)](h_0)>w+d,
\]
which is $Y_3>w+d$, hence $\Psi_{\rm R1}<0$.  The actual-reach handoff and
geometric contradiction occur only after this scalar inequality and are not
used here.
\end{proof}

\begin{theorem}[Pure S2-R2]
\label{thm:s2-univ-r2}
For every $(r,a,d)\in\mathcal D_2$,
\[
 \Psi_{\rm R2}(r,a,d)<0.
\]
\end{theorem}

\begin{proof}
Put $Q=k/(2r)=Y_0$ and $X=r-d$.  The proof of
Lemma~\ref{lem:signed-ce2-two-threshold} uses only
$\Delta_R>0$ and $\Delta_L>0$ and gives the inclusive alternative
\[
 e(a)<Q\quad\text{or}\quad e(d)<Q.
\]
The scalar part of Proposition~\ref{prop:signed-ce2-one-gap}, before its
geometric contradiction, uses this alternative and the exact branches of
$T$ to prove
\[
 [T(1-d,\cdot)\mid T(1-m,\cdot)\mid T(1-a,\cdot)](X)>1-Q.
\]
No actual-reach or placement assertion occurs in that calculation.  Applying
the exact reflection duality three times shows that the negation of this
inequality is equivalent to the negation of
\[
 [T(1-a,\cdot)\mid T(1-m,\cdot)\mid T(1-d,\cdot)](Q)>1-X.
\]
The left side is $Y_3$ and $1-X=w+d$, so this is exactly
$\Psi_{\rm R2}<0$.  The threshold alternative is inclusive: the proof
requires at least one threshold and does not assert uniqueness.
\end{proof}

\subsection{The finite high-radial T3 endpoint problem}

The domain $\Omega_{\rm T3}$ includes only $c_1,c_5>1/2$.  The geometric
branch with $c_1\le1/2$ is the direct right-linear case handled before the
optimization bridge and by Lemma~\ref{thm:s2-geom-t3}.

\begin{theorem}[Bridge-realized S2-T3 scalar conclusion]
\label{thm:s2-source-t3}
Let $\xi\in\Omega_{\rm T3}$ be supplied by
Proposition~\ref{prop:s2-t3-source-bridge} from a support-isolated geometric
T3 configuration and classified by the authenticated $407X$ case analysis.
Then
\[
 \Psi_{\rm T3}(\xi)<0.
\]
\end{theorem}

\begin{proof}
Fix
$\xi=(r,a,d,\tau,\beta,a_1,a_5,c_1,c_5)\in\Omega_{\rm T3}$.
If $a_1+a_5>1$, Lemma~\ref{lem:s2-pure-cap} gives
\[
 F(c_5,a_5)+F(c_1,a_1)
 \le2-a_1-a_5<1.
\]
Assume $a_1+a_5\le1$.  The two contact predicates in the definition of the
hard cell assign a unique pair
$(\ell_5,\ell_1)\in\mathfrak B^2$.

For a hard cell, use the signed-center substitution
\[
 R=r,\qquad W=w,\qquad E=e,
 \qquad \alpha_C=a,\qquad\delta_C=d.
\]
For the T3 source put
\[
 D_T=\frac1z,
 \qquad R_T=\frac\tau z,
 \qquad \alpha_T=\beta,
 \qquad p_T=z-\beta.
\]
Then, directly from $z^2=1-\tau+\tau^2$,
\[
 R_T^2-D_TR_T+D_T^2=1,
 \qquad
 \alpha_T+p_T=z=\frac1{D_T},
\]
and the exact source quantities in Lemma~\ref{lem:407-four-label-loss} become
\[
 q_T=1-p_T=q,
 \qquad
 \frac{D_Tq_T}{R_T}=c_\star,
 \qquad
 q_T+\frac{1-R_T}{D_T}=u_\star.
\]
For the bridge-realized point, the predicates $\mathcal R_i$,
$\mathcal L_{\varepsilon,j}$, and $\mathcal K_h$ are the residual orders and
radial hit/miss orders supplied by the geometric classification.  Finally,
$\mathrm{Cell}_{\ell_5}(c_5,a_5)$ and
$\mathrm{Cell}_{\ell_1}(c_1,a_1)$ select the same explicit branch formulas
of $F$ as its four-branch table.

After these substitutions, the applicable row of the authenticated
$407X$ geometric calculation proves
\[
 F(c_5,a_5)+F(c_1,a_1)<1
\]
for the bridge-realized $\xi$.  Hence $\Psi_{\rm T3}(\xi)<0$.  This argument
does not assert that the $407X$ rows prove the same sign for every abstract
point of $\Omega_{\rm T3}$.
\end{proof}

\subsection{The adjacent exceptional triangle problems}

\begin{theorem}[Pure S2-SC]
\label{thm:s2-univ-sc}
The two inequalities in Problem~\ref{prob:s2-rescuer} hold on
$\Omega_{\rm SC}^{T}$ and $\Omega_{\rm SC}^{V}$, respectively.
\end{theorem}

\begin{proof}
For the first cell, $1-u=(1-x)\rho(\theta)$ and
$a=x\rho(\theta)$, so $\vartheta(a,u)=x$.  Since
$0<\rho(\theta)\le1$, also $a\le x$.  Put
\[
 \chi(y)=\frac{y(2-y)}{1+y}.
\]
On $[0,1/2]$ this function is increasing and
$\chi(1-M_c^{\rm sup})=c$.  The exact calculation in
Lemma~\ref{lem:short-t3-rescuer-profile} proves, on the displayed
$\theta$-interval and $x$-interval,
\[
 Q_\theta(x)=2x^2+(\theta-1)x+\theta\ge0.
\]
This is equivalent to $\chi(x)\le x+\theta=c$, and hence
$x\le 1-M_c^{\rm sup}$.  Therefore
$\max\{a,\vartheta(a,u)\}\le 1-M_c^{\rm sup}$.

For the second cell put
\[
 G(t,c)=\Delta(t)-1-\frac{3t}{2}+c(t+1).
\]
The two defining equations and the midpoint inequality give $a\le G(t,c)$.
For $t\ge1$ and $c\le1/2$,
\[
 \partial_tG
 \le\frac{2t+1}{2\Delta(t)}-1<0,
\]
so
\[
 a\le G(t,c)\le L_{\rm sc}(c):=\sqrt3-\frac52+2c.
\]
The exact polynomial calculation in
Lemma~\ref{lem:short-vd1-rescuer-profile} gives
$2L_{\rm sc}(c)\le 1-M_c^{\rm sup}$ on $0\le c\le1/2$.
Writing $\epsilon=1-u$ gives
\[
 \epsilon=\epsilon_0+\frac at,
 \qquad
 \epsilon_0=\frac{2t+1-c(t+1)-\Delta(t)}t>0,
 \qquad
 \epsilon_0+\frac{G(t,c)}t=\frac12.
\]
Monotonicity in $a$ yields
\[
 \vartheta(a,u)
 =\frac a{a+\epsilon}
 \le\frac{G(t,c)}{G(t,c)+1/2}
 \le2L_{\rm sc}(c)
 \le 1-M_c^{\rm sup}.
\]
Thus both terms in the maximum are at most $1-M_c^{\rm sup}$, which is
exactly $\Psi_{\rm SC}\le0$.
\end{proof}

\subsection{The Vd radial problems}

\begin{lemma}[Pure adjacent quarter margin]
\label{lem:s2-pure-adjacent-quarter}
On either adjacent Vd cell,
\[
 d<\frac{A_L}{4}.
\]
\end{lemma}

\begin{proof}
By the definition of $\mathcal R_{[k/w,r+a]}$, either
$A_L=w-a$ or $A_L=1-p_0$.
If $A_L=w-a$, then
\[
 4d+a\le4(a+d)<\frac{2w}{3}<w,
\]
so $4d<A_L$.

Suppose $A_L=1-p_0$.  If also $A_R=1-q_0$, then
$A_R+A_L<1/2$ gives
$p_0+q_0>3/2$, whereas
$p_0^2+p_0q_0+q_0^2\le1$ gives
$p_0+q_0\le2/\sqrt3<3/2$.  Hence $A_R=r-d$ and $q_0\ge k/r$.
The endpoint-distance inequality gives $1-p_0>q_0/2$, so
\[
 A_L>\frac{k}{2r}>\frac\eta{2r}.
\]
Since $A_R+A_L<1/2$, also
\[
 d>J(r):=r-\frac12+\frac\eta{2r}.
\]
The exact derivative calculation from the proof of
Lemma~\ref{lem:signed-vd-adjacent-placement} gives $J'(r)>0$ and
$J(3/8)=1/24$.  Since $d<1/24$, it follows that $r<3/8$.  Then
\[
 (2-3r)^2-e^2=(3-8r)w>0,
\]
which implies $\eta/r>1/3$.  Hence $A_L>1/6$, while $4d<1/6$, and again
$4d<A_L$.
\end{proof}

\begin{theorem}[Pure S2-VD-adjacent]
\label{thm:s2-univ-vd-adjacent}
Both inequalities in Problem~\ref{prob:s2-vd-adjacent} hold on their stated
cells.
\end{theorem}

\begin{proof}
Lemma~\ref{lem:s2-pure-adjacent-quarter} gives the first strict term.  On the
positive-support cell, $t>0$ and
$\Delta_t<t+1$, so
\[
 u_+
 <1-q_1-\frac{p_1}{t}
 \le1-A_L
 <1-d.
\]
Thus $u_+-1+d<0$, and both entries of
$\Psi_{\rm VD}^{\rm adj,+}$ are negative.  On the no-support cell the
objective is only $d-A_L/4<0$.
\end{proof}

\begin{theorem}[Pure S2-VD-nonadjacent]
\label{thm:s2-univ-vd-nonadjacent}
The inequality in Problem~\ref{prob:s2-vd-nonadjacent} holds on every stated
cell.
\end{theorem}

\begin{proof}
Put $B_R=1-A_R$ and $B_L=1-A_L$.  The two explicit center-distance constraints in
Definition~\ref{def:s2-vd-nonadjacent-domain} imply
\[
 w+d\le M_0(x),
 \qquad
 r+a\le M_0(y).
\]
Here $M_0$ is the canonical zero-radial envelope from the common transfer
calculus, not a new local transfer map.
If $B_R=q_0$ and $p_0<x$, then $B_L=p_0$ and $A_R+A_L<1/2$ would give
$p_0+q_0>3/2$, contradicting the endpoint-distance inequality.  Hence
$p_0\ge x$, and
$q_0\le M_0(p_0)\le M_0(x)$.  The other possibility is
$B_R=w+d\le M_0(x)$.  Therefore
\[
 A_R\ge1-M_0(x)>\frac x2>\sigma.
\]
The reflected argument gives $A_L>y/2>\sigma$.  Hence
$\sigma-A_R<0$ and $\sigma-A_L<0$.

For $j=2,4$ one has $d_j\in\{d,a\}<\sigma$.  For $j=3$, at least one of
$a/r$ and $d/w$ is at most $a+d=\sigma$; otherwise multiplication by
$r,w$ and addition would give $a+d>a+d$.  Thus in every case
$d_j\le\sigma<m=\min\{A_R,A_L\}$.
Finally, $t>0$ gives $\Delta_t<t+1$, and therefore
\[
 C_v
 <1-\frac{p_v+tq_v}{t+1}
 \le1-\min\{p_v,q_v\}
 \le1-m.
\]
All four entries in $\Psi_{\rm VD}^{\rm non}$ are strictly negative.
\end{proof}

\begin{corollary}[Source-only scalar collection]
\label{cor:s2-source-registry}
Problems S2-E1, S2-E2, S2-R1, S2-R2, S2-SC, and S2-VD hold with the
universal quantifiers in their definitions.  The S2-T3 conclusion holds for
the bridge-realized geometric points stated in
Theorem~\ref{thm:s2-source-t3}; no universal T3 assertion is made here.
\end{corollary}

\begin{proof}
Combine Theorems~\ref{thm:s2-univ-e1}, \ref{thm:s2-univ-e2},
\ref{thm:s2-univ-r1}, \ref{thm:s2-univ-r2},
\ref{thm:s2-source-t3}, \ref{thm:s2-univ-sc},
\ref{thm:s2-univ-vd-adjacent}, and
\ref{thm:s2-univ-vd-nonadjacent}.
\end{proof}
